\cotefile{75}
\scanpage{18}
\textcircled{2} Donnée d'un $G$-torseur \inserted{$T$} $\unsure{\longleftrightarrow|\longrightarrow}$ donnée d'un ensemble $F$ muni d'une orbite de $G$ opérant sur $\operatorname{Bij}(E,F)$
% slanted
\margin{i.e. $F$ muni d'une $(G,E)$ structure}
[\unsure{quel|quand|où} $E$, \deleted{\unsure{opérant|opère} \unsure{fid} sur $E$, $\operatorname{Bij}(E,F)$ \unsure{est}}\inserted{$G$-ens. fidèle, est} donné une fois pour \unsure{ttes|toutes}] C'est un cas particulier de \textcircled{1}, car $G=\operatorname{Aut}_{C}(E)$, $C$ étant l'espèce de $G$-structure de type $(G,E)$.

\textcircled{3} Soit $C$ une espèce de structure (sur un ens de base, disons), $E$ une structure de type $C$, \inserted{$\unsure{C|G}=\operatorname{Aut}_{C}(E)$.} Alors, pour un ens $F$, la donnée d'une structure de type $C$ \uline{isom} à $E$ \uline{équivaut} à la donnée d'une $(G,E)$-structure.

De cette façon, \textcircled{1} devient cas particulier de \textcircled{2}.

\textcircled{4} Enveloppe vectorielle d'un polygône combinatoire.

On considère le polygône \deleted{\unsure{stan|st}} combinatoire \unsure{standard|standart} \deleted{$\unsure{[1,n]|[1,n[}$\unsure{,}} $P_n$, sa réalisation \inserted{(rectiligne)} régulière standard $i\mathrel{\unsure{\longrightarrow|\longmapsto}}\exp\frac{2i\pi}{n}=\zeta^i$ (il faut prouver que c'est bien une réalisation rectiligne \inserted{si $n\geq3$}\ldots !)

On note que $D_n=\operatorname{Aut}(P_n)$ opère sur \deleted{$\unsure{\mathbb{C}|E}$} \unsure{le|l'espace} plan euclidien \inserted{$E=$} $\mathbb{C}$ (muni de $|z|$) par \deleted{$\unsure{\tau_i|\sigma_i},\unsure{\sigma|\tau}\;\unsure{\longmapsto|\longrightarrow}$}
\[
\left\{\begin{aligned}
\tau_i &\longrightarrow \text{homothétie par }\zeta^i\\
\sigma &\longmapsto z\longmapsto\bar z
\end{aligned}\right.
\]
% archive sheet number: 17
\scanpage{19}
La donnée d'un polygone combinatoire $P$ à $n$ côtés équivaut à la donnée d'un $D_n$-torseur $T$ \inserted{$(P={}^{T}P_n)$}, et à $T$ on associe $\,{}^{\unsure{T|T'}}E$

L'inclusion $S(P_n)\hookrightarrow E$ (\unsure{étant|c'est} une $\deleted{\unsure{P}}D_n$-application) définit une inclusion
\[
 S(P)={}^{T}S(P_n)\hookrightarrow{}^{T}E=\operatorname{Env}(P)
\]

\noindent\mbox{-{}-{}-{}-{}-}
\[
 (e,G)\mathrel{\overset{\approx}{\unsure{-|\rightarrow}}}\underline{\operatorname{Tors}}(G)
\]

% slanted
\margin{DEA}
\noindent\(\left\{\begin{minipage}{0.93\linewidth}
\uline{NB}
\diagram{\begin{tikzcd}[ampersand replacement=\&,column sep=small]
\underline{\operatorname{Hom}}(\underline{\operatorname{Tors}}(G),C)
 \arrow[r,"\approx"] \&
 \begin{gathered}\text{catégor. }\deleted{\underline{\operatorname{Hom}}(G,\mathcal E)}
 \inserted{G\text{-}(\mathcal E)}\\
 \text{des objets }X\text{ de }C\\
 \text{munis d'une opération}\\
 \text{à }\deleted{\text{droite}}\inserted{\text{gauche}}\text{ de }G\\
 G\longrightarrow\operatorname{Aut}(X)^{\unsure{\mathcal E|\epsilon}}
 \end{gathered}\\
 \deleted{\unsure{T_X|P_X}\;\unsure{d'|d}} \& {}\\
 F\arrow[r,mapsto] \& F(G_d)\\
 (P\longmapsto{}^{P}X) \& X\arrow[l,mapsto]
\end{tikzcd}}

Soit $C$ catégorie, $E\in\operatorname{Ob} C$, $C_E$ \unsure{ss/cat.|s/cat.} des objets de $C$ isom. à $E$, avec iso, $G\unsure{=|:}\operatorname{Aut}_C(E)$
\[
 C_E\xrightarrow{\approx}\operatorname{Tors}(G)
\]
\diagram{\scriptsize\begin{tikzcd}[ampersand replacement=\&,row sep=large,column sep=small]
 \underline{\operatorname{Hom}}(\operatorname{Tors}(G),\mathcal E)
 \arrow[r,"\approx","{\unsure{\alpha\,\simeq|\alpha\,?}}"']
 \arrow[d,"\simeq"] \&
 \deleted{\underline{\operatorname{Hom}}(G,\mathcal E)}\inserted{G\text{-}(\mathcal E)}
 \arrow[dl,dashed,bend left=13,"\approx"']\\
 \underline{\operatorname{Hom}}(C_E,\mathcal E)
 \arrow[ur,leftrightarrow,bend left=7,"\approx",
 "{\Phi\mapsto\Phi_{\deleted{\unsure{E|0}}}(E)}"']
 \& X\longmapsto(F\longmapsto{}^{\operatorname{Isom}(E,F)}X\unsure{)|,})\\
 {}\arrow[ur,shorten <=1em,shorten >=1.2em] \& {}
\end{tikzcd}}

Donc : la catégorie des \unsure{foncteurs} de $C_E$ ds $\mathcal E$ est \unsure{équivalente} (de façon can. à isom unique près) à celle des $G$-objets de $\mathcal E$
\end{minipage}\right.\)

% slanted
\margin{\uline{cf} $C_0$ groupoïde \unsure{connexe}, $\unsure{E_0|E}\in\operatorname{Ob} C_0$, $G=\operatorname{Aut}(E_0)$ \par
$\underline{\operatorname{Hom}}(C_0,\mathcal E)\longrightarrow G\text{-}(\mathcal E)$ \par
$\Phi\longmapsto\Phi(E)$ \unsure{d'équivalence|équivalence} de catégories.}
% archive sheet number: 18
\scanpage{20}
Supposons qu'à tt polygône comb. de \unsure{rg|rang} $n$ (disons) $P$ on veuille faire correspondre un $\mathbb{R}$-vectoriel \deleted{\unsure{(un $\mathbb{R}$-|sur $\mathbb{R}$)}} (disons), ``de façon compatible avec le transport de structure'', \deleted{\unsure{Alors|Al}} soit $E(P)$. Alors $E(P_n)\inserted{=E}$ sera un vectoriel, avec $D_n=\operatorname{Aut}(P_n)$ opérant (à g.). Mais la donnée de $E$ avec cette opération permet de récupérer $E(P_{\deleted{\unsure{n}}})$ pour tt $P$, comme étant \deleted{$\unsure{E^T|E'}$} $\,{}^TE$, où $T=\operatorname{Isom}(P_n,P)$.

\textcircled{5} \deleted{\unsure{Changmt de|Chang$^{\rm mt}$ de}}\inserted{Extension} du groupe d'opérateurs

$G\mathrel{\unsure{-|\longrightarrow}}G'$, $E$ $G$-ens \inserted{(à \deleted{\unsure{g}}dr)}\unsure{,} on définit
\[
 {}^{\deleted{\unsure{G'|T}}}E^{G'}=\deleted{\unsure{G'\times_G|G'\times E}}E\times_G G'
 \quad\text{comme $G'$-ens. à g.}
\]

\uline{NB} Si $F$ un $G'$-ens. à \unsure{drte|dr.}, donc
\[
\left\{\begin{aligned}
 &G\text{ ens à droite, on a}\\
 &(E\times_G G')\times_{G'}F\simeq E\times_G(G'\times_{G'}F)\simeq E\times_G F\\
 &\deleted{F\times_{G'}(G'\times_G E)\simeq(F\times_{G'}\unsure{G'|(\,')})\times_G E\simeq F\times_G E}\quad\text{i.e.}
\end{aligned}\right.
\]
\[
 \deleted{F\times_{G'}E'\simeq F\times_G E}\qquad E'\times_{G'}F\simeq E\times_G F
\]

Si $E\inserted{=T}$ torseur \deleted{$\unsure{T|?}$}, $E'\inserted{=T'}$ aussi, donc on a
\[
 {}^{T}F^{\deleted{\unsure{T|T'}}}\simeq{}^{T'}F^{\deleted{\unsure{T'|T}}}
\]

i.e. Pour tordre $F$ par \inserted{$G$-}torseur $T$, quand $G$ opère
% archive sheet number: 19
\scanpage{21}
\uline{via} opération d'un groupe $G'$, il suffit de tordre par le torseur ${}^{\deleted{\unsure{G'|G}}}T^{G'}=T'$. (Ex : \deleted{\unsure{Si|Soit}} $F$ muni de \inserted{quelque} structure, \deleted{\unsure{disons,|disons, à}} on prend $G'=\operatorname{Aut}(F)^{\deleted{\unsure{\bullet|\ast}}}$\ldots{} et $F^T$ n'est autre \deleted{\illegible{1}} que la structure déduite du $G'$-torseur $T'$\ldots)

\textcircled{6} Opérations \inserted{de produit} sur les torseurs sur groupe \uline{commutatif}. Op. à g ou à droite pareil

$X\times_G Y\times_G Z$ ou \unsure{en|enc} $X\wedge Y\wedge Z$ produit associatif et commutatif (à isom can. près) $\unsure{\mathbb{1}|\underline{1}}=G_s$ \unsure{pour|joue} rôle d'une unité
\margin{$X\times_GY\simeq(X\times Y)\times_{G\times G}G$}
\[
 (X\times Y\times\unsure{-|\cdots})\times_{G^{\inserted{n}}\deleted{\unsure{\times\cdots|n\times}}}G
\]

\uline{Prop} Un produit de torseurs est un \uline{torseur}.

\deleted{\unsure{Soit|Soit $Y$}}

\uline{Prop} $\forall$ torseur $X$, $\exists$ torseur $\deleted{\unsure{X}}Y$ et isom $X\wedge\deleted{\unsure{X}}Y\simeq\deleted{\unsure{G_s|G}}\underline{\mathbb{1}}$. Pour un autre $Y'$, $\exists$ isom unique $Y\simeq Y'$ tel que.
\[
 Y\simeq X\times_G(G\overset{g\mapsto-g}{\mathrel{\unsure{\longrightarrow|-}}}G)
\]

\textcircled{7} \uline{Ex} \deleted{Si \unsure{on a une famille|on a une} d'ens. à $2$ \unsure{élémts|éléments}} \inserted{Opération $X\wedge Y$ entre ens. à $2$ élémts.}

\deleted{\unsure{Les|Ces} torseurs}

\deleted{\uline{Ex} Soit $B$ ens. fini avec rel. d'équiv. $R$ à orbites de cardinal $2$, $B\longrightarrow I=B/R$ Considérons $\bigwedge_{i\in I}^{\unsure{\ast|\bullet}}B_i$, un élément}
% archive sheet number: 20
\scanpage{22}
\textcircled{8} Digression sur structures cubiques.

\[
 A\text{ anneau avec }2\ne0,\qquad n\in\mathbb{N}^{*}
\]

$\unsure{\mathcal{C}|C}$ \inserted{cat} modules libres \inserted{$E$} sur $A$, avec \deleted{\unsure{famille|fam. de}} ens. de \inserted{$n$} sous-modules libres $E_i$ ($i\in I$, $\operatorname{card} I=n$) de somme directe $E$, et chaque $E_i$ muni d'une bi-base $\{e_i,f_i\}$ ($e_i\unsure{+}f_i=0$)

[ou encore, $E$ muni d'une partie $B$ de \inserted{card. $2n$} stable par $x\longmapsto-x$, et telle que un \textendash{} donc tt \textendash{} syst. de représentants de $B/\pm1$ \unsure{dans|de} $B$ soit une base]

\diagram{\begin{tikzcd}[ampersand replacement=\&,row sep=small]
 {}\arrow[d,"\unsure{\sim|?}"']\\ {}
\end{tikzcd}}
$\mathcal{C}'$ cat. des ens \inserted{$B$} de cardinal $2n$ munis d'une rel. d'équivalence \deleted{\unsure{bi à|bi}} à orbites de card. $2$, ou encore muni d'une permutation $\sigma$ involut\unsure{ive|ive.} ($\sigma^{\unsure{2|\prime}}=\mathrm{id}$) sans pt fixe.

\diagram{\begin{tikzcd}[ampersand replacement=\&,column sep=large]
 C\arrow[r]\& C' \quad\text{évident}\\
 \mathcal{C}'\arrow[r]\& C\quad\text{: si }I=\deleted{\unsure{E}}B/\sigma,\ \text{on}
\end{tikzcd}}
considère $E=\bigoplus_{i\in I}E_i$, où $E_i=B_i\wedge_{\{\pm1\}}\deleted{\unsure{B|A}}D$ où $D$ est le $A$-module $A_{\unsure{g|G}}$ sur lequel $\pm1$ opère par \unsure{homothétie|homothetie}
% archive sheet number: 21
\scanpage{23}
\uline{Contour \inserted{(comb.)} connexe} \inserted{$C$} définit

\textcircled{1} L'ens $\omega=\omega(C)$ à deux \unsure{éléments|élémts} de \unsure{ses|ces} deux orientations.

À partir de là, on peut déjà écrire (à isom can. près) le groupe des translations du contour, \unsure{par}
\[
 G^+=\operatorname{Aut}^+(C)\xleftarrow{\sim}\omega\wedge_{\{\pm1\}}\mathbb Z_n
\]
(où $\{\pm1\}$ opère sur $\mathbb Z_n$ par \unsure{homothéties|homotheties})

\textcircled{2} Un torseur sous $G^+$

\deleted{\uline{\unsure{Exercice|Exercices}}} Et \unsure{inversmt|inversemt|inversement}, le couple formé par $(\omega,T)$, où $\omega$ un \deleted{\unsure{tors|ens}} à $2$ él et $T$ un torseur sous $G^+\mathrel{\unsure{=|\simeq|\leftrightarrows}}\omega\wedge_{\{\pm1\unsure{,}\}}\mathbb Z_n$, définit un
% archive sheet number: 22
\scanpage{24}
contour

\uline{\unsure{Exos|Exo}} (DEA) Soit $G=Z.\Gamma$ produit $1/2$ direct de $\deleted{\unsure{I|\Gamma}}Z$ et $\deleted{\unsure{Z|I}}\Gamma$ ($\Gamma$ opérant sur $Z$). Montrer que la \unsure{donnée|donné} d'un torseur \inserted{$T$} sous $G$ équivaut \unsure{à|:} celle de

a) d'un torseur $\omega$ sous $\Gamma$

b) d'un torseur $T_0$ sous $\,{}^\omega Z=\omega\wedge_\Gamma Z$

\figure{\begin{tikzpicture}[x=1cm,y=1cm]
\draw[line width=.48pt,line cap=round] (0,0) .. controls (.48,-.65) and (1.35,-1.84) .. (2.08,-2.48);
\end{tikzpicture}}
% archive sheet number: 23
\scanpage{25}
\uline{Orientations}

a) Vectoriel réel $E$ de dim finie $n$. La donnée de $E$ équivaut à celle du torseur à dr.
\[
 P=\operatorname{Isom}(\mathbb R^n,E)\simeq\operatorname{bases}(E)
\]
sous
\[
 G=\operatorname{Gl}(n,\mathbb R).
\]

On a un hom
\diagram{\begin{tikzcd}[ampersand replacement=\&,column sep=small,row sep=small]
 G=\operatorname{Gl}(n,\deleted{\unsure{R|k}}\mathbb R)\arrow[rr,"s"]\arrow[dr,"\det"']
 \& {} \& \{\pm1\} \& u\longmapsto\deleted{\unsure{\det|d}}\,\operatorname{sgn}\det u\\
 {} \& \mathbb R^{\unsure{\ast|\times}}\arrow[ur,"\operatorname{sgn}"'] \& {}
 \& \text{noyau }\operatorname{Gl}^+(n,\mathbb R)\,\unsure{,}
\end{tikzcd}}
et
\[
 \omega_E=P\times_G\{\pm1,s\}\simeq P/\operatorname{Gl}^+(n,R)
\]
est l'ens des orientations de $E$.

b) Module \inserted{\unsure{quadratique|quadratiques}} libre \deleted{\unsure{quadr.|point}} \inserted{$E$} sur un anneau comm. $k$ \unsure{(où|où} $2$ est inversible [$\exists$ base orthonormale, $\deleted{\unsure{Q(e_i,e|Q(e_i,e_j)}} e_i^2=1$, $e_i e_j=0$ si $i\ne j$]

La donnée de $E$ de rang $n$ équivaut à celle du tors à droite
\[
 P=\operatorname{Isom}_{k\text{-quadr.}}(k^n,E)\simeq\text{bases orthon.}(E)
\]
sous
\[
 G=\operatorname{O}(n,k)=\operatorname{Aut}_{k\text{-quadr}}(k^n,Q_n)\subset\operatorname{Gl}(n\unsure{,|.}k)
\]

\unsure{Or|On} ici
\[
 G=\operatorname{O}(n,k)\xrightarrow{\det=s}\{\pm1\}\qquad(\text{noyau }\operatorname{SO}(n,k))
\]
et l'ens des orientations de $E$ est (par définition)
\[
 \omega_{\unsure{E|F}}=P\times_G(\pm1,s)\simeq P/\operatorname{SO}(n,k).
\]

\uline{Compatibilité} \uline{entre a) et b)} Résulte de la transitivité de l'opération d'ext. du groupe d'opérateurs dans un torseur, et du diagramme comm.
\diagram{\begin{tikzcd}[ampersand replacement=\&,column sep=large]
 \operatorname{O}(n,\unsure{k|\mathbb R})\arrow[r,hook,"\mathrm{incl}"]\arrow[rr,bend right=19,"\det"] \& \operatorname{Gl}(n,k)\arrow[r,"\operatorname{sgn}\det"] \& \{\pm1\}
\end{tikzcd}}
% archive sheet number: 24
